\section*{Chapter \ref{ch:pl:python-at-scale} --- Python at Scale}

\begin{solution}{exo:pl:python-at-scale:1}
$960 \times 2\times10^8$~ns $= 192$~s, about three minutes; $145$~ns: 29~s; $14$~ns: 2.7~s.
\end{solution}

\begin{solution}{exo:pl:python-at-scale:2}
Each field access on a mapped record builds a record object and then a scalar object, with their reference counting and
destruction; the loop over lists reads integers that already exist. The arithmetic is the same; the objects are not.
\end{solution}

\begin{solution}{exo:pl:python-at-scale:3}
About $3 \times 61 = 183$~MB as objects and $3 \times 1.0 = 3$~MB as categoricals.
\end{solution}

\begin{solution}{exo:pl:python-at-scale:4}
With $c_0 = 0$ and $c_t = \sum_{s\le t} x_s$ (a cumulative sum), the sum of the last $w$ values ending at $t$ is $c_t - c_{t-w}$ (with
$c_{t-w}$ taken as 0 when $t < w$). A chunked version carries the last $w-1$ values of the previous chunk (or, equivalently, the last
$w$ cumulative sums), so that the first $w-1$ windows of a chunk are complete.
\end{solution}

\begin{solution}{exo:pl:python-at-scale:5}
Two processes would each have their own interpreter and lock and run the two loops in parallel, in about the time of one. With a
64~MB input each, pickled through a pipe, each would first pay about 0.14~s of serialisation, more than the loop itself; through
shared memory or a mapped file, about 0.03~s, or nothing if the processes map the file themselves.
\end{solution}

\begin{solution}{exo:pl:python-at-scale:6}
$1.5\times10^9 / 16 \approx 94$ million events, about 0.03\% of a day of 311 billion messages: a large feed cannot be held whole in one
process; it must be streamed.
\end{solution}

\begin{solution}{exo:pl:python-at-scale:7}
The differences are of the order of $10^{-13}$ for every chunk size: rounding, since the chunked version's cumulative sums start from
each chunk's carried tail rather than from the first value; they are exact in the mathematical sense.
\end{solution}

\begin{solution}{exo:pl:python-at-scale:8}
A pure-Python loop holds the interpreter lock: eight threads run it one at a time, in about the time of one thread over the whole day,
plus switching. Use eight processes, each mapping its part of the file, or rewrite the loop in arrays; and if the feature is recursive
across the day, carry its state across the parts rather than restarting it at each boundary.
\end{solution}

\begin{solution}{pb:pl:python-at-scale:1}
\begin{enumerate}
\item It builds and destroys record and scalar objects, looks up names, dispatches every operation on types, and counts references,
for each of the ten or so bytecode operations of an iteration.
\item About 960 (loop over records), 145 (loop over lists), 14 (array expression) and 15 (chunked).
\item 70 times over the records, 11 times over the lists.
\item The filter performs the recursion's multiplications and additions in the same order as the loop.
\item Loops whose state changes shape with each event (a book, a state machine with many branches); they go to compiled code
(chapter~9).
\item About 61, 1, 8 and 16 bytes.
\item Fifty distinct symbols stored once, and a one-byte code per row: \omterm{def:pl:columnar-formats:encodings}{dictionary encoding} in memory.
\item About 16 bytes an event: the signed quantities and the output, each 8 bytes, plus short-lived temporaries.
\item About 94 million.
\item The largest physical memory the process used at any moment; a job is killed at its peak, not at its average.
\item The filter's last output starts the next chunk's recursion (\cref{prop:pl:python-at-scale:exact}).
\item Chunks of 100\,000 events, with a peak of 2.5~MB; chunks of 10\,000 are within ten per cent at 0.31~MB.
\item The interpreter's overhead is paid once per chunk, two thousand times for two million events.
\item On a Python loop, nothing (twice the work in twice the time); on a numpy sort, most of a second core, since the sort releases the
lock (0.094~s for two against 0.086~s for one).
\item About 0.14~s pickled, 0.03~s through shared memory, copy included.
\item \textbf{Named result.} The array expression is 70 times faster than a Python loop over the mapped records and 11 times faster
than one over lists (14 against 960 and 145 nanoseconds an event); streamed in chunks of 10\,000 events it takes about the same time
(0.025~s against 0.028~s whole, 0.023~s at the fastest chunk of 100\,000) with a peak of 0.31~MB instead of 32~MB, so a 1.5~GB budget holds about 94 million events whole and any history streamed.
\item $250 \times 2\times10^8 = 5\times10^{10}$ events: about 13 hours by the loop over records, 2 hours over lists, 11 minutes by
arrays.
\item Eight processes, each mapping its share of the days' files and streaming them in chunks, each returning a small result; no data
through pipes.
\item It would let threads run Python loops in parallel, so the threaded version of exercise~8 would scale; it would not change the
interpreter's overhead per event, the memory of objects, or the value of streaming.
\item Keeping the loops out of the interpreter and the whole history out of memory.
\end{enumerate}
\end{solution}

\begin{solution}{iq:pl:python-at-scale:1}
Each iteration runs in the interpreter, building and checking objects for every element: hundreds of nanoseconds against a few
for the arithmetic. Express it as array operations (cumulative sums, filters, \texttt{searchsorted}), or compile the loop.
\iqlookfor{\omterm{def:pl:python-at-scale:overhead}{Interpreter overhead} per element; the array idioms.}
\end{solution}

\begin{solution}{iq:pl:python-at-scale:2}
As a linear filter (one call in compiled code), streamed over the data in chunks from a \omterm{def:pl:a-tick-store-on-open-formats:mmap}{memory-mapped file}, with the last output of
each chunk as the next chunk's initial condition; memory stays at one chunk.
\iqlookfor{Recursion as a filter; chunking with carried state.}
\end{solution}

\begin{solution}{iq:pl:python-at-scale:3}
A lock that lets one thread at a time run Python bytecode. It does not matter for input and output, or for compiled code that releases
it (much of numpy), and it is sidestepped by processes; a free-threaded build of recent CPython versions removes it.
\iqlookfor{Concurrency against parallelism; what releases the lock.}
\end{solution}

\begin{solution}{iq:pl:python-at-scale:4}
Object columns (strings as Python objects: convert to categoricals or Arrow strings), wide types (float64 or int64 where 32 or 8 bits
suffice), copies held by intermediate variables, and whether the whole history needs to be in memory at all.
\iqlookfor{Representations first, then streaming.}
\end{solution}

\begin{solution}{iq:pl:python-at-scale:5}
Partition the work by days (or instruments) into a queue; run about 30 worker processes, each mapping or scanning its partition from
the store and streaming it in chunks with carried state, writing results as files; size chunks so that 30 workers stay well within
128~GB; measure time and peak memory on one day before launching all.
\iqlookfor{Processes over threads, data-local work, a measured memory budget.}
\end{solution}
